\chapter{Hito 8: Construcción de las Tablas de Hechos en DuckDB}

\section{Objetivo}

El objetivo de este hito es construir las tablas de hechos del Data Warehouse macrofinanciero utilizando DuckDB.

Las tablas de hechos almacenan las medidas numéricas del modelo dimensional. En este proyecto construiremos dos hechos principales:

\begin{itemize}
\item \texttt{fact_tipo_cambio}
\item \texttt{fact_balance_mensual}
\end{itemize}

Al finalizar este hito el estudiante será capaz de:

\begin{itemize}
\item Crear tablas de hechos en DuckDB.
\item Cargar hechos desde la capa \texttt{staging}.
\item Resolver claves subrogadas utilizando las dimensiones.
\item Calcular medidas derivadas.
\item Validar el grano de cada tabla de hechos.
\item Preparar la información para construir data marts.
\end{itemize}

\section{Resultado esperado}

Al finalizar este hito deberán existir:

\begin{verbatim}
core.fact_tipo_cambio

core.fact_balance_mensual
\end{verbatim}

Estas tablas deberán estar conectadas lógicamente con:

\begin{verbatim}
core.dim_tiempo

core.dim_moneda

core.dim_entidad_financiera
\end{verbatim}

\section{Arquitectura del hito}

El flujo será:

\begin{verbatim}
staging.tipo_cambio_raw
|
V
core.fact_tipo_cambio

staging.balance_bancario_raw
|
V
core.fact_balance_mensual
\end{verbatim}

\section{Paso 1: Crear script DDL de hechos}

Crear:

\begin{verbatim}
sql/05_facts_ddl.sql
\end{verbatim}

\section{Paso 2: Crear tabla fact_tipo_cambio}

Agregar:

\begin{lstlisting}[language=SQL]
CREATE TABLE IF NOT EXISTS core.fact_tipo_cambio
(
sk_tiempo INTEGER,

```
sk_moneda INTEGER,

tasa_compra DOUBLE,

tasa_venta DOUBLE,

tasa_promedio DOUBLE,

spread DOUBLE
```

);
\end{lstlisting}

\section{Paso 3: Crear tabla fact_balance_mensual}

Agregar:

\begin{lstlisting}[language=SQL]
CREATE TABLE IF NOT EXISTS core.fact_balance_mensual
(
sk_tiempo INTEGER,

```
sk_entidad INTEGER,

activos_totales DOUBLE,

cartera_creditos DOUBLE,

depositos DOUBLE,

patrimonio DOUBLE,

resultado_neto DOUBLE,

depositos_me DOUBLE,

cartera_me DOUBLE
```

);
\end{lstlisting}

\section{Nota sobre restricciones}

DuckDB permite definir llaves primarias y restricciones, pero para este primer proyecto mantendremos el diseño simple.

La validación del grano se realizará mediante consultas de control.

\section{Paso 4: Crear script de carga de hechos}

Crear:

\begin{verbatim}
sql/06_load_facts.sql
\end{verbatim}

\section{Paso 5: Limpiar hechos}

Agregar:

\begin{lstlisting}[language=SQL]
DELETE FROM core.fact_tipo_cambio;

DELETE FROM core.fact_balance_mensual;
\end{lstlisting}

\section{Paso 6: Cargar fact_tipo_cambio}

Agregar:

\begin{lstlisting}[language=SQL]
INSERT INTO core.fact_tipo_cambio
(
sk_tiempo,
sk_moneda,
tasa_compra,
tasa_venta,
tasa_promedio,
spread
)

SELECT

```
t.sk_tiempo,

m.sk_moneda,

CAST(r.tasa_compra AS DOUBLE)
    AS tasa_compra,

CAST(r.tasa_venta AS DOUBLE)
    AS tasa_venta,

(
    CAST(r.tasa_compra AS DOUBLE)
    +
    CAST(r.tasa_venta AS DOUBLE)
) / 2
    AS tasa_promedio,

(
    CAST(r.tasa_venta AS DOUBLE)
    -
    CAST(r.tasa_compra AS DOUBLE)
)
    AS spread
```

FROM staging.tipo_cambio_raw r

JOIN core.dim_tiempo t
ON t.fecha = CAST(r.fecha AS DATE)

JOIN core.dim_moneda m
ON m.moneda_cod = r.moneda

WHERE r.fecha IS NOT NULL
AND r.moneda IS NOT NULL
AND CAST(r.tasa_compra AS DOUBLE) > 0
AND CAST(r.tasa_venta AS DOUBLE) > 0;
\end{lstlisting}

\section{Explicación}

En esta transformación:

\begin{itemize}
\item Se convierte la fecha desde texto hacia \texttt{DATE}.
\item Se convierten las tasas desde texto hacia \texttt{DOUBLE}.
\item Se resuelve la clave \texttt{sk_tiempo}.
\item Se resuelve la clave \texttt{sk_moneda}.
\item Se calcula \texttt{tasa_promedio}.
\item Se calcula \texttt{spread}.
\end{itemize}

\section{Paso 7: Cargar fact_balance_mensual}

Agregar:

\begin{lstlisting}[language=SQL]
INSERT INTO core.fact_balance_mensual
(
sk_tiempo,
sk_entidad,
activos_totales,
cartera_creditos,
depositos,
patrimonio,
resultado_neto,
depositos_me,
cartera_me
)

SELECT

```
t.sk_tiempo,

e.sk_entidad,

CAST(r.activos_totales AS DOUBLE)
    AS activos_totales,

CAST(r.cartera_creditos AS DOUBLE)
    AS cartera_creditos,

CAST(r.depositos AS DOUBLE)
    AS depositos,

CAST(r.patrimonio AS DOUBLE)
    AS patrimonio,

CAST(r.resultado_neto AS DOUBLE)
    AS resultado_neto,

CAST(r.depositos_me AS DOUBLE)
    AS depositos_me,

CAST(r.cartera_me AS DOUBLE)
    AS cartera_me
```

FROM staging.balance_bancario_raw r

JOIN core.dim_tiempo t
ON t.fecha =
last_day(
CAST(
r.periodo || '01'
AS DATE
)
)

JOIN core.dim_entidad_financiera e
ON e.entidad_cod = r.entidad_cod

WHERE r.periodo IS NOT NULL
AND r.entidad_cod IS NOT NULL;
\end{lstlisting}

\section{Advertencia sobre fechas mensuales}

Si el campo \texttt{periodo} viene con formato \texttt{YYYYMM}, DuckDB puede requerir una conversión explícita.

Si el bloque anterior falla, utilizar esta versión alternativa:

\begin{lstlisting}[language=SQL]
STRPTIME(
r.periodo || '01',
'%Y%m%d'
)::DATE
\end{lstlisting}

La unión con \texttt{dim_tiempo} quedaría:

\begin{lstlisting}[language=SQL]
JOIN core.dim_tiempo t
ON t.fecha =
last_day(
STRPTIME(
r.periodo || '01',
'%Y%m%d'
)::DATE
)
\end{lstlisting}

\section{Paso 8: Crear script Python para cargar hechos}

Crear:

\begin{verbatim}
etl/load_facts.py
\end{verbatim}

Contenido:

\begin{lstlisting}[language=Python]
from db import get_connection

def run_sql_file(path, conn):

```
with open(
    path,
    "r",
    encoding="utf-8"
) as file:

    sql = file.read()

conn.execute(sql)
```

def main():

```
conn = get_connection()

run_sql_file(
    "sql/05_facts_ddl.sql",
    conn
)

run_sql_file(
    "sql/06_load_facts.sql",
    conn
)

conn.close()

print(
    "Hechos cargados correctamente."
)
```

if **name** == "**main**":
main()
\end{lstlisting}

\section{Paso 9: Ejecutar carga de hechos}

Desde la terminal:

\begin{lstlisting}[language=bash]
python etl/load_facts.py
\end{lstlisting}

Salida esperada:

\begin{verbatim}
Hechos cargados correctamente.
\end{verbatim}

\section{Paso 10: Crear consulta de monitoreo}

Crear:

\begin{verbatim}
sql/check_facts.sql
\end{verbatim}

Contenido:

\begin{lstlisting}[language=SQL]
SELECT
'fact_tipo_cambio' AS tabla,
COUNT(*) AS total_filas
FROM core.fact_tipo_cambio

UNION ALL

SELECT
'fact_balance_mensual' AS tabla,
COUNT(*) AS total_filas
FROM core.fact_balance_mensual;
\end{lstlisting}

\section{Paso 11: Crear script para revisar hechos}

Crear:

\begin{verbatim}
etl/check_facts.py
\end{verbatim}

Contenido:

\begin{lstlisting}[language=Python]
from db import get_connection

conn = get_connection()

with open(
"sql/check_facts.sql",
"r",
encoding="utf-8"
) as file:

```
sql = file.read()
```

result = conn.execute(sql).fetchdf()

print(result)

conn.close()
\end{lstlisting}

Ejecutar:

\begin{lstlisting}[language=bash]
python etl/check_facts.py
\end{lstlisting}

\section{Paso 12: Inspeccionar fact_tipo_cambio}

Crear:

\begin{verbatim}
sql/inspect_fact_tipo_cambio.sql
\end{verbatim}

Contenido:

\begin{lstlisting}[language=SQL]
SELECT

```
t.fecha,

m.moneda_cod,

f.tasa_compra,

f.tasa_venta,

f.tasa_promedio,

f.spread
```

FROM core.fact_tipo_cambio f

JOIN core.dim_tiempo t
ON t.sk_tiempo = f.sk_tiempo

JOIN core.dim_moneda m
ON m.sk_moneda = f.sk_moneda

ORDER BY
t.fecha,
m.moneda_cod

LIMIT 10;
\end{lstlisting}

\section{Paso 13: Inspeccionar fact_balance_mensual}

Crear:

\begin{verbatim}
sql/inspect_fact_balance_mensual.sql
\end{verbatim}

Contenido:

\begin{lstlisting}[language=SQL]
SELECT

```
t.fecha,

e.entidad_cod,

e.entidad_nombre,

f.activos_totales,

f.depositos,

f.patrimonio,

f.resultado_neto
```

FROM core.fact_balance_mensual f

JOIN core.dim_tiempo t
ON t.sk_tiempo = f.sk_tiempo

JOIN core.dim_entidad_financiera e
ON e.sk_entidad = f.sk_entidad

ORDER BY
t.fecha,
e.entidad_cod

LIMIT 10;
\end{lstlisting}

\section{Paso 14: Validar el grano de tipo de cambio}

Crear:

\begin{verbatim}
sql/check_grain_fx.sql
\end{verbatim}

Contenido:

\begin{lstlisting}[language=SQL]
SELECT

```
sk_tiempo,

sk_moneda,

COUNT(*) AS total
```

FROM core.fact_tipo_cambio

GROUP BY
sk_tiempo,
sk_moneda

HAVING COUNT(*) > 1;
\end{lstlisting}

Resultado esperado:

\begin{verbatim}
0 filas
\end{verbatim}

\section{Paso 15: Validar el grano de balance mensual}

Crear:

\begin{verbatim}
sql/check_grain_balance.sql
\end{verbatim}

Contenido:

\begin{lstlisting}[language=SQL]
SELECT

```
sk_tiempo,

sk_entidad,

COUNT(*) AS total
```

FROM core.fact_balance_mensual

GROUP BY
sk_tiempo,
sk_entidad

HAVING COUNT(*) > 1;
\end{lstlisting}

Resultado esperado:

\begin{verbatim}
0 filas
\end{verbatim}

\section{Paso 16: Validar medidas derivadas}

Crear:

\begin{verbatim}
sql/check_fx_measures.sql
\end{verbatim}

Contenido:

\begin{lstlisting}[language=SQL]
SELECT

```
tasa_compra,

tasa_venta,

tasa_promedio,

spread,

(tasa_compra + tasa_venta) / 2
    AS tasa_promedio_recalculada,

tasa_venta - tasa_compra
    AS spread_recalculado
```

FROM core.fact_tipo_cambio

LIMIT 10;
\end{lstlisting}

\section{Paso 17: Validar reglas de negocio}

Crear:

\begin{verbatim}
sql/check_balance_business_rules.sql
\end{verbatim}

Contenido:

\begin{lstlisting}[language=SQL]
SELECT *
FROM core.fact_balance_mensual
WHERE depositos_me > depositos;

SELECT *
FROM core.fact_balance_mensual
WHERE cartera_me > cartera_creditos;
\end{lstlisting}

Resultado esperado:

\begin{verbatim}
0 filas
\end{verbatim}

\section{Paso 18: Integrar al pipeline}

Actualizar:

\begin{verbatim}
etl/run_pipeline.py
\end{verbatim}

Contenido recomendado:

\begin{lstlisting}[language=Python]
import sys

import subprocess

steps = [
"etl/load_staging.py",
"etl/load_dimensions.py",
"etl/load_facts.py"
]

for step in steps:

```
print(
    f"Ejecutando {step}..."
)

subprocess.run(
    [
        sys.executable,
        step
    ],
    check=True
)
```

print(
"Pipeline ejecutado correctamente."
)
\end{lstlisting}

\section{Buenas prácticas}

Durante este hito seguiremos estas reglas:

\begin{enumerate}
\item Los hechos deben respetar el grano definido.
\item Los hechos deben usar claves subrogadas.
\item Las medidas derivadas deben calcularse de forma reproducible.
\item Las validaciones deben ejecutarse después de cada carga.
\item El pipeline debe poder ejecutarse de inicio a fin.
\end{enumerate}

\section{Checklist de validación}

\begin{itemize}
\item[$\Box$] Tabla \texttt{fact_tipo_cambio} creada.
\item[$\Box$] Tabla \texttt{fact_balance_mensual} creada.
\item[$\Box$] Hechos cargados.
\item[$\Box$] Conteos revisados.
\item[$\Box$] Grano de tipo de cambio validado.
\item[$\Box$] Grano de balance mensual validado.
\item[$\Box$] Medidas derivadas verificadas.
\item[$\Box$] Pipeline actualizado.
\end{itemize}

\section{Entregable}

Al finalizar este hito el estudiante deberá disponer de:

\begin{itemize}
\item Tablas de hechos en \texttt{core}.
\item Scripts SQL de creación y carga.
\item Consultas de inspección.
\item Validaciones de grano.
\item Modelo dimensional funcional.
\end{itemize}

\section{Próximo hito}

En el Hito 9 construiremos la capa \texttt{mart}.

Aprenderemos a:

\begin{itemize}
\item Crear vistas analíticas.
\item Calcular indicadores macrofinancieros.
\item Construir ROA, ROE, dolarización y HHI.
\item Preparar datos listos para Streamlit.
\end{itemize}
